% Options for packages loaded elsewhere
\PassOptionsToPackage{unicode}{hyperref}
\PassOptionsToPackage{hyphens}{url}
\documentclass[
]{article}
\usepackage{xcolor}
\usepackage[margin=1in]{geometry}
\usepackage{amsmath,amssymb}
\setcounter{secnumdepth}{-\maxdimen} % remove section numbering
\usepackage{iftex}
\ifPDFTeX
  \usepackage[T1]{fontenc}
  \usepackage[utf8]{inputenc}
  \usepackage{textcomp} % provide euro and other symbols
\else % if luatex or xetex
  \usepackage{unicode-math} % this also loads fontspec
  \defaultfontfeatures{Scale=MatchLowercase}
  \defaultfontfeatures[\rmfamily]{Ligatures=TeX,Scale=1}
\fi
\usepackage{lmodern}
\ifPDFTeX\else
  % xetex/luatex font selection
\fi
% Use upquote if available, for straight quotes in verbatim environments
\IfFileExists{upquote.sty}{\usepackage{upquote}}{}
\IfFileExists{microtype.sty}{% use microtype if available
  \usepackage[]{microtype}
  \UseMicrotypeSet[protrusion]{basicmath} % disable protrusion for tt fonts
}{}
\makeatletter
\@ifundefined{KOMAClassName}{% if non-KOMA class
  \IfFileExists{parskip.sty}{%
    \usepackage{parskip}
  }{% else
    \setlength{\parindent}{0pt}
    \setlength{\parskip}{6pt plus 2pt minus 1pt}}
}{% if KOMA class
  \KOMAoptions{parskip=half}}
\makeatother
\usepackage{color}
\usepackage{fancyvrb}
\newcommand{\VerbBar}{|}
\newcommand{\VERB}{\Verb[commandchars=\\\{\}]}
\DefineVerbatimEnvironment{Highlighting}{Verbatim}{commandchars=\\\{\}}
% Add ',fontsize=\small' for more characters per line
\usepackage{framed}
\definecolor{shadecolor}{RGB}{248,248,248}
\newenvironment{Shaded}{\begin{snugshade}}{\end{snugshade}}
\newcommand{\AlertTok}[1]{\textcolor[rgb]{0.94,0.16,0.16}{#1}}
\newcommand{\AnnotationTok}[1]{\textcolor[rgb]{0.56,0.35,0.01}{\textbf{\textit{#1}}}}
\newcommand{\AttributeTok}[1]{\textcolor[rgb]{0.13,0.29,0.53}{#1}}
\newcommand{\BaseNTok}[1]{\textcolor[rgb]{0.00,0.00,0.81}{#1}}
\newcommand{\BuiltInTok}[1]{#1}
\newcommand{\CharTok}[1]{\textcolor[rgb]{0.31,0.60,0.02}{#1}}
\newcommand{\CommentTok}[1]{\textcolor[rgb]{0.56,0.35,0.01}{\textit{#1}}}
\newcommand{\CommentVarTok}[1]{\textcolor[rgb]{0.56,0.35,0.01}{\textbf{\textit{#1}}}}
\newcommand{\ConstantTok}[1]{\textcolor[rgb]{0.56,0.35,0.01}{#1}}
\newcommand{\ControlFlowTok}[1]{\textcolor[rgb]{0.13,0.29,0.53}{\textbf{#1}}}
\newcommand{\DataTypeTok}[1]{\textcolor[rgb]{0.13,0.29,0.53}{#1}}
\newcommand{\DecValTok}[1]{\textcolor[rgb]{0.00,0.00,0.81}{#1}}
\newcommand{\DocumentationTok}[1]{\textcolor[rgb]{0.56,0.35,0.01}{\textbf{\textit{#1}}}}
\newcommand{\ErrorTok}[1]{\textcolor[rgb]{0.64,0.00,0.00}{\textbf{#1}}}
\newcommand{\ExtensionTok}[1]{#1}
\newcommand{\FloatTok}[1]{\textcolor[rgb]{0.00,0.00,0.81}{#1}}
\newcommand{\FunctionTok}[1]{\textcolor[rgb]{0.13,0.29,0.53}{\textbf{#1}}}
\newcommand{\ImportTok}[1]{#1}
\newcommand{\InformationTok}[1]{\textcolor[rgb]{0.56,0.35,0.01}{\textbf{\textit{#1}}}}
\newcommand{\KeywordTok}[1]{\textcolor[rgb]{0.13,0.29,0.53}{\textbf{#1}}}
\newcommand{\NormalTok}[1]{#1}
\newcommand{\OperatorTok}[1]{\textcolor[rgb]{0.81,0.36,0.00}{\textbf{#1}}}
\newcommand{\OtherTok}[1]{\textcolor[rgb]{0.56,0.35,0.01}{#1}}
\newcommand{\PreprocessorTok}[1]{\textcolor[rgb]{0.56,0.35,0.01}{\textit{#1}}}
\newcommand{\RegionMarkerTok}[1]{#1}
\newcommand{\SpecialCharTok}[1]{\textcolor[rgb]{0.81,0.36,0.00}{\textbf{#1}}}
\newcommand{\SpecialStringTok}[1]{\textcolor[rgb]{0.31,0.60,0.02}{#1}}
\newcommand{\StringTok}[1]{\textcolor[rgb]{0.31,0.60,0.02}{#1}}
\newcommand{\VariableTok}[1]{\textcolor[rgb]{0.00,0.00,0.00}{#1}}
\newcommand{\VerbatimStringTok}[1]{\textcolor[rgb]{0.31,0.60,0.02}{#1}}
\newcommand{\WarningTok}[1]{\textcolor[rgb]{0.56,0.35,0.01}{\textbf{\textit{#1}}}}
\usepackage{graphicx}
\makeatletter
\newsavebox\pandoc@box
\newcommand*\pandocbounded[1]{% scales image to fit in text height/width
  \sbox\pandoc@box{#1}%
  \Gscale@div\@tempa{\textheight}{\dimexpr\ht\pandoc@box+\dp\pandoc@box\relax}%
  \Gscale@div\@tempb{\linewidth}{\wd\pandoc@box}%
  \ifdim\@tempb\p@<\@tempa\p@\let\@tempa\@tempb\fi% select the smaller of both
  \ifdim\@tempa\p@<\p@\scalebox{\@tempa}{\usebox\pandoc@box}%
  \else\usebox{\pandoc@box}%
  \fi%
}
% Set default figure placement to htbp
\def\fps@figure{htbp}
\makeatother
\setlength{\emergencystretch}{3em} % prevent overfull lines
\providecommand{\tightlist}{%
  \setlength{\itemsep}{0pt}\setlength{\parskip}{0pt}}
\usepackage{bookmark}
\IfFileExists{xurl.sty}{\usepackage{xurl}}{} % add URL line breaks if available
\urlstyle{same}
\hypersetup{
  pdftitle={WAIC},
  hidelinks,
  pdfcreator={LaTeX via pandoc}}

\title{WAIC}
\author{}
\date{\vspace{-2.5em}2026-04-17}

\begin{document}
\maketitle

\subsection{Introduction}\label{introduction}

The widely applicable information criterion (WAIC), introduced by
Watanabe, estimates the expected out-of-sample prediction score from
posterior samples without holding out data. It is the Bayesian successor
to AIC and DIC: AIC requires a single point estimate (and assumes
regularity that fails for hierarchical models); DIC plug-ins a posterior
mean of the deviance and is unstable in mixed models. WAIC instead
averages the pointwise log-posterior predictive density across posterior
draws and penalises by an effective-parameter term, making it directly
applicable to any model with a well-defined pointwise likelihood.

\subsection{Prerequisites}\label{prerequisites}

A working understanding of information criteria, posterior predictive
distributions, and the relationship between in-sample fit and
out-of-sample prediction.

\subsection{Theory}\label{theory}

For data \(y_1, \dots, y_n\) and posterior draws \(\theta^{(s)}\), the
log pointwise predictive density is

\[\widehat{\mathrm{lpd}} = \sum_{i=1}^n \log \!\left(\frac{1}{S} \sum_{s=1}^S p(y_i \mid \theta^{(s)})\right).\]

The effective parameter count \(p_{\text{waic}}\) is the sum of
pointwise posterior variances of the log-likelihood. Then

\[\mathrm{WAIC} = -2\bigl(\widehat{\mathrm{lpd}} - p_{\text{waic}}\bigr).\]

Lower is better. Watanabe's asymptotic theory shows that for regular
models WAIC is asymptotically equivalent to leave-one-out
cross-validation. The standard error of a WAIC difference between two
models comes from the sample standard deviation of the pointwise
differences scaled by \(\sqrt n\), providing a principled threshold for
``meaningful'' differences.

\subsection{Assumptions}\label{assumptions}

A pointwise likelihood is well-defined (so each observation contributes
one term to the sum) and the posterior is well-mixed. WAIC inherits the
exchangeability assumption of standard cross-validation --- for
clustered or temporal data, replace it with leave-one-cluster-out or
leave-future-out variants.

\subsection{R Implementation}\label{r-implementation}

\begin{Shaded}
\begin{Highlighting}[]
\FunctionTok{library}\NormalTok{(brms); }\FunctionTok{library}\NormalTok{(loo)}

\FunctionTok{data}\NormalTok{(mtcars)}
\NormalTok{fit1 }\OtherTok{\textless{}{-}} \FunctionTok{brm}\NormalTok{(mpg }\SpecialCharTok{\textasciitilde{}}\NormalTok{ wt, }\AttributeTok{data =}\NormalTok{ mtcars, }\AttributeTok{chains =} \DecValTok{2}\NormalTok{, }\AttributeTok{iter =} \DecValTok{2000}\NormalTok{, }\AttributeTok{refresh =} \DecValTok{0}\NormalTok{)}
\NormalTok{fit2 }\OtherTok{\textless{}{-}} \FunctionTok{brm}\NormalTok{(mpg }\SpecialCharTok{\textasciitilde{}}\NormalTok{ wt }\SpecialCharTok{+}\NormalTok{ hp, }\AttributeTok{data =}\NormalTok{ mtcars, }\AttributeTok{chains =} \DecValTok{2}\NormalTok{, }\AttributeTok{iter =} \DecValTok{2000}\NormalTok{, }\AttributeTok{refresh =} \DecValTok{0}\NormalTok{)}

\NormalTok{w1 }\OtherTok{\textless{}{-}} \FunctionTok{waic}\NormalTok{(fit1)}
\NormalTok{w2 }\OtherTok{\textless{}{-}} \FunctionTok{waic}\NormalTok{(fit2)}

\FunctionTok{loo\_compare}\NormalTok{(w1, w2)}
\end{Highlighting}
\end{Shaded}

\subsection{Output \& Results}\label{output-results}

\texttt{waic()} returns the WAIC value, \(p_{\text{waic}}\), and a
vector of pointwise contributions. \texttt{loo\_compare()} aligns two
WAIC objects, reports their differences and standard errors, and orders
models from best (top) to worst (bottom).

\subsection{Interpretation}\label{interpretation}

A reporting sentence: ``WAIC favoured the two-predictor model with
\(\Delta\mathrm{WAIC} = -12 \pm 4\) (standard error of the difference);
under the rule-of-thumb that differences exceeding
\(2 \times \mathrm{SE}\) are meaningful, the data clearly support
including horsepower in addition to weight.'' Always quote both the
difference and its standard error --- a ``12-point improvement'' without
context is uninformative.

\subsection{Practical Tips}\label{practical-tips}

\begin{itemize}
\tightlist
\item
  Use PSIS-LOO from the \texttt{loo} package whenever possible; it is
  more numerically stable than WAIC and provides Pareto-\(k\)
  diagnostics that flag influential observations one at a time.
\item
  WAIC can be unstable when a single observation contributes a large
  fraction of the variance of the log-likelihood; if you see
  \(p_{\text{waic}}\) comparable to the number of observations, do not
  trust the criterion and switch to LOO.
\item
  Differences in WAIC are noisy; only invest belief in differences that
  exceed 2 standard errors.
\item
  WAIC is well-defined for non-nested models, unlike the
  likelihood-ratio test, which makes it natural for comparing models
  with different families or link functions.
\item
  For hierarchical models, choose the leave-out unit deliberately:
  leaving out a single observation may understate predictive uncertainty
  if all other observations from the same cluster remain in the training
  fold.
\end{itemize}

\subsection{R Packages Used}\label{r-packages-used}

\texttt{loo} for the canonical \texttt{waic()} and \texttt{loo()}
implementations and \texttt{loo\_compare()}, \texttt{brms} and
\texttt{rstanarm} for fits that expose pointwise log-likelihoods
automatically, and \texttt{posterior} for the underlying draw arrays.

\subsection{For Reviewers}\label{for-reviewers}

\textbf{What to look for in a paper using this method.}

\begin{itemize}
\tightlist
\item
  \textbf{Common misapplications.}
\item
  \textbf{Diagnostics that should be reported but often aren't.}
\item
  \textbf{Red flags in tables and figures.}
\item
  \textbf{What to verify.}
\item
  \textbf{What an adequate Methods paragraph must contain.}
\end{itemize}

\subsection{See also --- labs in this
chapter}\label{see-also-labs-in-this-chapter}

\begin{itemize}
\tightlist
\item
  \href{labs/lab-bayesian-thinking-control-vs-decision-error.Rmd}{Bayesian
  thinking; α control vs decision error}
\item
  \href{labs/lab-brms-stan-loo-hierarchical-models.Rmd}{brms / Stan,
  LOO, hierarchical models}
\end{itemize}

\end{document}
